% STUDENT-SAFE SOURCE: this file contains no answer labels or explanations.
% Generated by generate_core_syntax_bank.py. Edit the generator, not this file.
\documentclass[9pt]{extarticle}

% Generated by generate_core_syntax_bank.py. Edit the generator, not this file.
\usepackage[letterpaper,margin=0.55in,includehead,includefoot,
  headheight=15pt,headsep=6pt,footskip=18pt]{geometry}
\usepackage[T1]{fontenc}
\usepackage{amsmath,amssymb}
\usepackage{microtype}
\usepackage[dvipsnames]{xcolor}
\usepackage{multicol}
\usepackage{fancyhdr}
\usepackage[hidelinks]{hyperref}

\definecolor{CoreNavy}{HTML}{18364A}
\definecolor{CoreTeal}{HTML}{176B6B}
\definecolor{CoreTint}{HTML}{EDF4F8}

\newif\ifanswers
\ifdefined\MAKEANSWERKEY
  \answerstrue
\else
  \answersfalse
\fi

\setlength{\parindent}{0pt}
\setlength{\parskip}{0pt}
\setlength{\columnsep}{0.25in}
\setlength{\columnseprule}{0.25pt}
\setlength{\emergencystretch}{1.2em}
\raggedbottom

\newcommand{\currentcore}{Core Practice}
\newcommand{\currentform}{}
\pagestyle{fancy}
\fancyhf{}
\fancyhead[L]{\textbf{COMP 2300}}
\fancyhead[C]{\textbf{\ifanswers Instructor Key\else Read --- Parse --- Judge\fi}}
\fancyhead[R]{\currentform}
\fancyfoot[L]{\ifanswers Answers, diagnoses, and minimal repairs\else Read it; parse it; then judge it.\fi}
\fancyfoot[C]{\thepage}
\fancyfoot[R]{20 items}
\renewcommand{\headrulewidth}{0.5pt}
\renewcommand{\footrulewidth}{0.3pt}

\newcounter{tfitem}
\newcommand{\CoreForm}[5]{%
  \clearpage
  \renewcommand{\currentcore}{#1 Core}%
  \renewcommand{\currentform}{Form #2}%
  \setcounter{tfitem}{0}%
  {\color{CoreNavy}\Large\bfseries #1 Core: Reading \& Interpreting --- Form #2}\par
  \vspace{2pt}\color{CoreNavy}\rule{\linewidth}{0.8pt}\color{black}\par
  \vspace{3pt}
  \ifanswers\else
    \textbf{Name:} \rule{0.42\linewidth}{0.4pt}\hfill
    \textbf{Date:} \rule{0.22\linewidth}{0.4pt}\par
    \vspace{4pt}
  \fi
  \fcolorbox{CoreTeal}{CoreTint}{\parbox{0.965\linewidth}{%
    \ifanswers\scriptsize\else\footnotesize\fi
    \textbf{Given.} #3\par
    \vspace{2pt}
    \textbf{Say the symbols as words.} #4\par
    \vspace{2pt}
    \textbf{For every item:} (1) say the whole statement in words; (2) lightly
    box its main relation, connective, or claim and mark its inputs or scope;
    (3) check whether it is complete and well formed here; (4) only then apply
    the definitions and circle an answer. \textbf{Decision key: N} = not a
    truth-evaluable proposition as written (incomplete, undefined, free
    variable, type error, or arity error); \textbf{T/F} = a complete proposition
    that is true/false. \textbf{N is not F. Do not circle before you parse.}}}
  \vspace{5pt}
  \begin{multicols}{2}
  \raggedcolumns
  #5
  \end{multicols}
}

\newcommand{\StageHeader}[2]{%
  \par\noindent
  {\color{CoreNavy}\footnotesize\bfseries #1}\par\nopagebreak
  {\color{CoreTeal}\scriptsize\itshape Reading lens: #2}\par\nopagebreak
  \vspace{1.5pt}
}

\newcommand{\StudentItem}[1]{%
  \stepcounter{tfitem}%
  \par\noindent
  \hangindent=3.25em\hangafter=1
  \makebox[1.65em][r]{\textbf{\thetfitem.}}\hspace{0.25em}%
  \makebox[1.25em][c]{\textbf{T}}\hspace{0.10em}%
  \makebox[1.25em][c]{\textbf{F}}\hspace{0.10em}%
  \makebox[1.25em][c]{\textbf{N}}\hspace{0.30em}%
  #1%
  \par\vspace{5.5pt}
}

\newcommand{\KeyItem}[5]{%
  \stepcounter{tfitem}%
  \par\noindent
  \hangindent=3.25em\hangafter=1
  \makebox[1.65em][r]{\textbf{\thetfitem.}}\hspace{0.25em}%
  \makebox[1.25em][c]{\textbf{T}}\hspace{0.10em}%
  \makebox[1.25em][c]{\textbf{F}}\hspace{0.10em}%
  \makebox[1.25em][c]{\textbf{N}}\hspace{0.30em}%
  #5%
  \par\nobreak
  {\scriptsize\color{CoreTeal}\hspace*{3.25em}%
    \textbf{[#1 --- #2]} #3%
    \if\relax\detokenize{#4}\relax\else\
      \textit{Minimal repair:} #4%
    \fi
  }\par
}

\begin{document}
\ifanswers\small\else\fontsize{10.8}{13.0}\selectfont\fi

\CoreForm{Logic}{A}{$p$ and $r$ are true; $q$ is false. Variables range over $\mathbb Z$. $P(n)$ means \emph{$n$ is even}, $Q(n)$ means \emph{$n>0$}, and $R(a,b)$ means \emph{$a<b$}. $P,Q$ are unary predicates and $R$ is binary.}{$\neg p$: not $p$; $p\land q$: $p$ and $q$; $p\lor q$: $p$ or $q$; $p\to q$: if $p$, then $q$; $p\leftrightarrow q$: $p$ exactly when $q$; $\forall$: for every; $\exists$: there exists; $\therefore$: therefore.}{
\StageHeader{1. Build and read propositions}{Locate every operator and all of its operands before assigning a truth value.}
\StudentItem{$q$}
\StudentItem{$\neg q$}
\StudentItem{$p\land r$}
\StudentItem{$q\lor q$}
\StudentItem{$p\to q$}
\StageHeader{2. Interpret connectives and equivalences}{Find the main connective first; then evaluate from the inside out.}
\StudentItem{$p\leftrightarrow r$}
\StudentItem{$p\lor\neg$ is a tautology.}
\StudentItem{$p\to q\equiv$.}
\StudentItem{$\neg(p\lor q)\equiv\neg p\lor\neg q$.}
\StudentItem{A truth table for $p,q$ has $4$ rows.}
\columnbreak
\StageHeader{3. Check predicates, types, and scope}{Check predicate arity, object types, and whether every variable is bound.}
\StudentItem{$P$.}
\StudentItem{$R(1)$.}
\StudentItem{$\forall n\,P(n)$.}
\StudentItem{$\exists n\,[P(n)\land Q(n)]$.}
\StudentItem{$\neg\exists n\,P(n)\equiv\exists n\,\neg P(n)$.}
\StageHeader{4. Read quantified claims and arguments}{Read quantifiers left to right, mark their scope, and separate premises from conclusions.}
\StudentItem{$\forall a\,\exists\,R(a,b)$.}
\StudentItem{$P(x)$.}
\StudentItem{$r,\ p\to r\ \therefore p$ is valid.}
\StudentItem{$p\lor,\ \neg p\ \therefore q$ is valid.}
\StudentItem{$1<2$.}
}
\CoreForm{Logic}{B}{$p$ and $r$ are true; $q$ is false. Variables range over $\mathbb Z$. $P(n)$ means \emph{$n$ is even}, $Q(n)$ means \emph{$n>0$}, and $R(a,b)$ means \emph{$a<b$}. $P,Q$ are unary predicates and $R$ is binary.}{$\neg p$: not $p$; $p\land q$: $p$ and $q$; $p\lor q$: $p$ or $q$; $p\to q$: if $p$, then $q$; $p\leftrightarrow q$: $p$ exactly when $q$; $\forall$: for every; $\exists$: there exists; $\therefore$: therefore.}{
\StageHeader{1. Build and read propositions}{Locate every operator and all of its operands before assigning a truth value.}
\StudentItem{$p\land$}
\StudentItem{$\neg p$}
\StudentItem{$p\land q$}
\StudentItem{$\lor q$}
\StudentItem{$p\to$}
\StageHeader{2. Interpret connectives and equivalences}{Find the main connective first; then evaluate from the inside out.}
\StudentItem{$p\leftrightarrow q$}
\StudentItem{$p\lor\neg p$ is a tautology.}
\StudentItem{$p\to q\equiv\neg p\lor q$.}
\StudentItem{$\neg(p\land)\equiv\neg p$.}
\StudentItem{A truth table for $p,q$ has $8$ rows.}
\columnbreak
\StageHeader{3. Check predicates, types, and scope}{Check predicate arity, object types, and whether every variable is bound.}
\StudentItem{$P(2)$.}
\StudentItem{$R(1,2)$.}
\StudentItem{$\forall\,P(n)$.}
\StudentItem{$\exists n\,[P(n)\land\neg P(n)]$.}
\StudentItem{$\neg\forall\,P(n)$.}
\StageHeader{4. Read quantified claims and arguments}{Read quantifiers left to right, mark their scope, and separate premises from conclusions.}
\StudentItem{$\forall a\,\exists b\,R(a,b)$.}
\StudentItem{$\forall x\,[P(x)\lor\neg P(x)]$.}
\StudentItem{$p,\ p\to\ \therefore r$ is valid.}
\StudentItem{$p\to q,\ \neg q\ \therefore\neg p$ is valid.}
\StudentItem{$2<1$.}
}
\CoreForm{Logic}{C}{$p$ and $r$ are true; $q$ is false. Variables range over $\mathbb Z$. $P(n)$ means \emph{$n$ is even}, $Q(n)$ means \emph{$n>0$}, and $R(a,b)$ means \emph{$a<b$}. $P,Q$ are unary predicates and $R$ is binary.}{$\neg p$: not $p$; $p\land q$: $p$ and $q$; $p\lor q$: $p$ or $q$; $p\to q$: if $p$, then $q$; $p\leftrightarrow q$: $p$ exactly when $q$; $\forall$: for every; $\exists$: there exists; $\therefore$: therefore.}{
\StageHeader{1. Build and read propositions}{Locate every operator and all of its operands before assigning a truth value.}
\StudentItem{$p$}
\StudentItem{$\neg$}
\StudentItem{$p\,q$}
\StudentItem{$q\lor r$}
\StudentItem{$q\to p$}
\StageHeader{2. Interpret connectives and equivalences}{Find the main connective first; then evaluate from the inside out.}
\StudentItem{$p\leftrightarrow\land q$}
\StudentItem{$p\land\neg p$ is a tautology.}
\StudentItem{$p\to q\equiv p\land q$.}
\StudentItem{$\neg(p\land q)\equiv\neg p\lor\neg q$.}
\StudentItem{A truth table for $p,q$ has $2^{\Box}$ rows.}
\columnbreak
\StageHeader{3. Check predicates, types, and scope}{Check predicate arity, object types, and whether every variable is bound.}
\StudentItem{$P(3)$.}
\StudentItem{$R(2,1)$.}
\StudentItem{$\forall n\,[P(n)\lor\neg P(n)]$.}
\StudentItem{$\exists n$.}
\StudentItem{$\neg\forall n\,P(n)\equiv\exists n\,\neg P(n)$.}
\StageHeader{4. Read quantified claims and arguments}{Read quantifiers left to right, mark their scope, and separate premises from conclusions.}
\StudentItem{$\exists b\,\forall a\,R(a,b)$.}
\StudentItem{$\forall x\,[P(x)\land\neg P(x)]$.}
\StudentItem{$p,\ p\to r\ \therefore r$ is valid.}
\StudentItem{$p\lor q,\ \neg p\ \therefore\neg q$ is valid.}
\StudentItem{$1+2$.}
}
\CoreForm{Arguments}{A}{$p,q,r$ denote propositions. A displayed argument lists premises before $\therefore$ and one conclusion after it. $\Gamma\models C$ means every valuation satisfying all premises in $\Gamma$ also satisfies $C$.}{$\neg$: not; $\land$: and; $\lor$: inclusive or; $p\to q$: if $p$, then $q$; $\Gamma\models C$: $C$ follows from premises $\Gamma$; $\therefore$: therefore.}{
\StageHeader{1. Recognize valid argument forms}{Separate the premises from the conclusion, then name the proposed inference pattern.}
\StudentItem{$p,\ p\to q\ \therefore$ is valid.}
\StudentItem{$p\to q,\ \neg q\ \therefore\neg p$ is valid.}
\StudentItem{$p\lor,\ \neg p\ \therefore q$ is valid.}
\StudentItem{$p\to q,\ r\to q\ \therefore p\to r$ is valid.}
\StudentItem{$p,\ q\ \therefore p\land q$ is valid.}
\StageHeader{2. Separate validity, soundness, and consistency}{Ask whether the premises could all be true while the conclusion is false.}
\StudentItem{$p\lor q\ \therefore p$ is valid.}
\StudentItem{Soundness $=$ valid $\land$.}
\StudentItem{A countermodel makes one premise false and the conclusion true.}
\StudentItem{An invalid argument can have a true conclusion.}
\StudentItem{$p$ and $\neg$ are inconsistent.}
\columnbreak
\StageHeader{3. Read assumptions, flaws, and conditionals}{Identify exactly what the conclusion needs and what each premise actually supplies.}
\StudentItem{A necessary assumption is merely an unrelated true statement.}
\StudentItem{$p\therefore$ because $p$.}
\StudentItem{A correlation alone does not establish which variable causes the other.}
\StudentItem{sample $\subseteq$.}
\StudentItem{In $p\to q$, $q$ is sufficient for $p$.}
\StageHeader{4. Test what the information supports}{Use only the stated information; distinguish support from proof and relevance from truth.}
\StudentItem{A strengthening answer must independently prove the conclusion.}
\StudentItem{alternative $\lor$.}
\StudentItem{A must-be-true answer follows from all the supplied facts.}
\StudentItem{A point-at-issue answer states something one speaker accepts and the other rejects.}
\StudentItem{Parallel reasoning requires the same nouns in both arguments.}
}
\CoreForm{Arguments}{B}{$p,q,r$ denote propositions. A displayed argument lists premises before $\therefore$ and one conclusion after it. $\Gamma\models C$ means every valuation satisfying all premises in $\Gamma$ also satisfies $C$.}{$\neg$: not; $\land$: and; $\lor$: inclusive or; $p\to q$: if $p$, then $q$; $\Gamma\models C$: $C$ follows from premises $\Gamma$; $\therefore$: therefore.}{
\StageHeader{1. Recognize valid argument forms}{Separate the premises from the conclusion, then name the proposed inference pattern.}
\StudentItem{$p,\ p\to q\ \therefore q$ is valid.}
\StudentItem{$p\to q,\ \neg p\ \therefore\neg q$ is valid.}
\StudentItem{$p\lor q,\ \neg p\ \therefore q$ is valid.}
\StudentItem{$p\to q,\ q\to\ \therefore p\to r$ is valid.}
\StudentItem{$p\ \therefore p\land q$ is valid.}
\StageHeader{2. Separate validity, soundness, and consistency}{Ask whether the premises could all be true while the conclusion is false.}
\StudentItem{$p\ \therefore\lor q$ is valid.}
\StudentItem{A sound argument is valid and has true premises.}
\StudentItem{$\Gamma\models$.}
\StudentItem{A true conclusion makes every argument for it valid.}
\StudentItem{$p$ and $\neg p$ are inconsistent.}
\columnbreak
\StageHeader{3. Read assumptions, flaws, and conditionals}{Identify exactly what the conclusion needs and what each premise actually supplies.}
\StudentItem{premise $\because$.}
\StudentItem{A circular argument uses its conclusion as support for itself.}
\StudentItem{A correlation by itself proves a causal mechanism.}
\StudentItem{A representative sample is relevant to a population generalization.}
\StudentItem{In $p\to$, the consequent is necessary for $p$.}
\StageHeader{4. Test what the information supports}{Use only the stated information; distinguish support from proof and relevance from truth.}
\StudentItem{premises $\uparrow$ conclusion $\because$.}
\StudentItem{A weakening answer may expose an alternative explanation.}
\StudentItem{A statement that is merely compatible with the facts must be true.}
\StudentItem{Different examples guarantee that two speakers disagree.}
\StudentItem{$p+q\ \therefore r$.}
}
\CoreForm{Arguments}{C}{$p,q,r$ denote propositions. A displayed argument lists premises before $\therefore$ and one conclusion after it. $\Gamma\models C$ means every valuation satisfying all premises in $\Gamma$ also satisfies $C$.}{$\neg$: not; $\land$: and; $\lor$: inclusive or; $p\to q$: if $p$, then $q$; $\Gamma\models C$: $C$ follows from premises $\Gamma$; $\therefore$: therefore.}{
\StageHeader{1. Recognize valid argument forms}{Separate the premises from the conclusion, then name the proposed inference pattern.}
\StudentItem{$q,\ p\to q\ \therefore p$ is valid.}
\StudentItem{$p\to,\ \neg q\ \therefore\neg p$ is valid.}
\StudentItem{$p\lor q,\ p\ \therefore\neg q$ is valid.}
\StudentItem{$p\to q,\ q\to r\ \therefore p\to r$ is valid.}
\StudentItem{$p,\ q\ \therefore p\land$ is valid.}
\StageHeader{2. Separate validity, soundness, and consistency}{Ask whether the premises could all be true while the conclusion is false.}
\StudentItem{$p\ \therefore p\lor q$ is valid.}
\StudentItem{A valid argument must have true premises.}
\StudentItem{A countermodel to validity makes every premise true and the conclusion false.}
\StudentItem{true $\therefore$.}
\StudentItem{$p$ and $p\lor q$ are inconsistent.}
\columnbreak
\StageHeader{3. Read assumptions, flaws, and conditionals}{Identify exactly what the conclusion needs and what each premise actually supplies.}
\StudentItem{A necessary assumption must hold for the stated support to work.}
\StudentItem{An argument becomes noncircular whenever its conclusion is true.}
\StudentItem{cause$(p,q,$.}
\StudentItem{Sample size and selection method never affect a generalization.}
\StudentItem{In $p\to q$, $q$ is necessary for $p$.}
\StageHeader{4. Test what the information supports}{Use only the stated information; distinguish support from proof and relevance from truth.}
\StudentItem{A strengthening answer adds support to the conclusion relative to the premises.}
\StudentItem{A weakening answer must establish the exact negation of the conclusion.}
\StudentItem{facts $\models$.}
\StudentItem{speaker A $\leftrightarrow$.}
\StudentItem{Parallel arguments can share a logical form while discussing different topics.}
}
\CoreForm{Sets}{A}{$U=\{1,2,3,4\}$, $A=\{1,2\}$, $B=\{2,3\}$, and $C=A$. Complements are relative to $U$. Let $f:A\to U$ satisfy $f(1)=2$, $f(2)=4$, and let $R=\{(1,1),(2,2)\}$ be a relation on $A$.}{$x\in A$: $x$ is an element of $A$; $A\subseteq B$: every $A$-element is in $B$; $\cup$: in either set; $\cap$: in both; $A\setminus B$: in $A$ but not $B$; $A^c$: outside $A$; $\mathcal P(A)$: the set of all subsets of $A$.}{
\StageHeader{1. Read membership, inclusion, and size}{Name the object on each side and distinguish membership from subset inclusion.}
\StudentItem{$3\in A$.}
\StudentItem{$A\subseteq$.}
\StudentItem{$A=\{1,2$.}
\StudentItem{$\varnothing\in A$.}
\StudentItem{$|A|=3$.}
\StageHeader{2. Interpret set operations}{Translate each operation into an element-by-element condition before computing.}
\StudentItem{$A\cup B=\{1,2,3\}$.}
\StudentItem{$A\cap\in B$.}
\StudentItem{$A\setminus B=\{2\}$.}
\StudentItem{$A^c=\{3,4\}$.}
\StudentItem{$(A\cup B)^c=A^c\cap B^c$.}
\columnbreak
\StageHeader{3. Read products and relations}{Check the type of every object: element, ordered pair, set, or relation.}
\StudentItem{$\mathcal P()=A$.}
\StudentItem{$\varnothing\in\mathcal P(A)$.}
\StudentItem{$(1,2)=(2,1)$.}
\StudentItem{$A\times\in B$.}
\StudentItem{$R$ is reflexive on $A$.}
\StageHeader{4. Read functions, variables, and types}{Bind every variable and distinguish a function value from a claim about the function.}
\StudentItem{$(1,2)\in R\Rightarrow(2,)\in R$.}
\StudentItem{$f(2)=2$.}
\StudentItem{$f$ is surjective onto $U$.}
\StudentItem{$\forall x\in A,\ x\in U$.}
\StudentItem{$A\cap B$.}
}
\CoreForm{Sets}{B}{$U=\{1,2,3,4\}$, $A=\{1,2\}$, $B=\{2,3\}$, and $C=A$. Complements are relative to $U$. Let $f:A\to U$ satisfy $f(1)=2$, $f(2)=4$, and let $R=\{(1,1),(2,2)\}$ be a relation on $A$.}{$x\in A$: $x$ is an element of $A$; $A\subseteq B$: every $A$-element is in $B$; $\cup$: in either set; $\cap$: in both; $A\setminus B$: in $A$ but not $B$; $A^c$: outside $A$; $\mathcal P(A)$: the set of all subsets of $A$.}{
\StageHeader{1. Read membership, inclusion, and size}{Name the object on each side and distinguish membership from subset inclusion.}
\StudentItem{$1\in$.}
\StudentItem{$A\subseteq U$.}
\StudentItem{$A=C$.}
\StudentItem{$\varnothing\subseteq$.}
\StudentItem{$|A|=$.}
\StageHeader{2. Interpret set operations}{Translate each operation into an element-by-element condition before computing.}
\StudentItem{$A\cup B=\{2\}$.}
\StudentItem{$A\cap B=\{2\}$.}
\StudentItem{$A\setminus$.}
\StudentItem{$A^c=\{1,2\}$.}
\StudentItem{$(A\cup B)^c=A^c\cup B^c$.}
\columnbreak
\StageHeader{3. Read products and relations}{Check the type of every object: element, ordered pair, set, or relation.}
\StudentItem{$|\mathcal P(A)|=4$.}
\StudentItem{$\{3\}\in\mathcal P(A)$.}
\StudentItem{$(1,)\in A\times B$.}
\StudentItem{$|A\times B|=4$.}
\StudentItem{$R$ is irreflexive on $A$.}
\StageHeader{4. Read functions, variables, and types}{Bind every variable and distinguish a function value from a claim about the function.}
\StudentItem{$R$ is symmetric on $A$.}
\StudentItem{$f(1,2)=4$.}
\StudentItem{$f(A)=2\land$.}
\StudentItem{$\forall x\in U,\ x\in A$.}
\StudentItem{$2\in A\cap B$.}
}
\CoreForm{Sets}{C}{$U=\{1,2,3,4\}$, $A=\{1,2\}$, $B=\{2,3\}$, and $C=A$. Complements are relative to $U$. Let $f:A\to U$ satisfy $f(1)=2$, $f(2)=4$, and let $R=\{(1,1),(2,2)\}$ be a relation on $A$.}{$x\in A$: $x$ is an element of $A$; $A\subseteq B$: every $A$-element is in $B$; $\cup$: in either set; $\cap$: in both; $A\setminus B$: in $A$ but not $B$; $A^c$: outside $A$; $\mathcal P(A)$: the set of all subsets of $A$.}{
\StageHeader{1. Read membership, inclusion, and size}{Name the object on each side and distinguish membership from subset inclusion.}
\StudentItem{$1\in A$.}
\StudentItem{$B\subseteq A$.}
\StudentItem{$A=B$.}
\StudentItem{$\varnothing\subseteq A$.}
\StudentItem{$|A|=2$.}
\StageHeader{2. Interpret set operations}{Translate each operation into an element-by-element condition before computing.}
\StudentItem{$A\cup=\{1,2,3\}$.}
\StudentItem{$A\cap B=\{1,2,3\}$.}
\StudentItem{$A\setminus B=\{1\}$.}
\StudentItem{$(A\cup)^c$.}
\StudentItem{$(A\cup B)^c=A^c\cap$.}
\columnbreak
\StageHeader{3. Read products and relations}{Check the type of every object: element, ordered pair, set, or relation.}
\StudentItem{$|\mathcal P(A)|=2$.}
\StudentItem{$\mathcal P(A)\in$.}
\StudentItem{$(1,2)\ne(2,1)$.}
\StudentItem{$|A\times B|=2$.}
\StudentItem{$1R$.}
\StageHeader{4. Read functions, variables, and types}{Bind every variable and distinguish a function value from a claim about the function.}
\StudentItem{$\{(1,2)\}$ is symmetric on $A$.}
\StudentItem{$f(1)=2$.}
\StudentItem{$f$ is injective.}
\StudentItem{$x\in A$.}
\StudentItem{$1\in A\cap B$.}
}
\CoreForm{Proofs}{A}{Variables range over $\mathbb Z$ when bound. $P$ and $Q$ denote propositions, and $P(n)$ denotes a proposition about an integer $n$. A formula with an unbound variable is not a complete proposition.}{$\forall$: for every; $\exists$: there exists; $P\to Q$: if $P$, then $Q$; $P\leftrightarrow Q$: both directions; a \emph{witness} proves an existential; a \emph{counterexample} disproves a universal.}{
\StageHeader{1. Identify the proof strategy}{State what may be assumed and what must be shown for the named strategy.}
\StudentItem{A direct proof of $P\to Q$ may assume $P$ and derive $Q$.}
\StudentItem{$\neg Q\to\neg$.}
\StudentItem{To prove $P$ by contradiction, one may assume $\neg P$ and derive $\bot$.}
\StudentItem{A proof by cases is complete after one convenient case.}
\StudentItem{$P\leftrightarrow$.}
\StageHeader{2. Read quantifiers, witnesses, and counterexamples}{Match universal claims with arbitrary objects and existential claims with witnesses.}
\StudentItem{A proof of a universal statement may choose only an object with a convenient extra property.}
\StudentItem{To prove $\exists n\,P(n)$, one must prove $P(n)$ for every integer.}
\StudentItem{From $\forall n$, infer $P(7)$.}
\StudentItem{From $\exists n\,P(n)$, one may introduce a fresh symbol $c$ with $P(c)$.}
\StudentItem{One integer $c$ with $P(c)$ disproves $\forall n\,P(n)$.}
\columnbreak
\StageHeader{3. Expand definitions inside proofs}{Expand the relevant definition before deciding whether a displayed line proves anything.}
\StudentItem{$9=2(4)$ proves that $9$ is even.}
\StudentItem{$2k+1$.}
\StudentItem{$\frac{3}{4}$ is rational.}
\StudentItem{$6\mid20$.}
\StudentItem{even $+$ odd $=$.}
\StageHeader{4. Audit proof language and logical structure}{Separate established facts from goals, conjectures, and incomplete proof language.}
\StudentItem{The product of two odd integers is odd.}
\StudentItem{A conjecture is already established by proof.}
\StudentItem{In $P\to Q$, $Q$ is necessary for $P$.}
\StudentItem{$P\because P\therefore$.}
\StudentItem{$n^2\ge0$.}
}
\CoreForm{Proofs}{B}{Variables range over $\mathbb Z$ when bound. $P$ and $Q$ denote propositions, and $P(n)$ denotes a proposition about an integer $n$. A formula with an unbound variable is not a complete proposition.}{$\forall$: for every; $\exists$: there exists; $P\to Q$: if $P$, then $Q$; $P\leftrightarrow Q$: both directions; a \emph{witness} proves an existential; a \emph{counterexample} disproves a universal.}{
\StageHeader{1. Identify the proof strategy}{State what may be assumed and what must be shown for the named strategy.}
\StudentItem{A direct proof of $P\to Q$ may assume $Q$ and derive $P$.}
\StudentItem{Proving $\neg Q\to\neg P$ proves $P\to Q$.}
\StudentItem{To prove $P$ by contradiction, assume $P$ and stop.}
\StudentItem{case 1 $\lor$.}
\StudentItem{To prove $P\leftrightarrow Q$, it is enough to prove $P\to Q$ and $Q\to P$.}
\StageHeader{2. Read quantifiers, witnesses, and counterexamples}{Match universal claims with arbitrary objects and existential claims with witnesses.}
\StudentItem{Let $n=$ be arbitrary.}
\StudentItem{$\exists\,P(n)$.}
\StudentItem{From $\forall n\,P(n)$, one may infer $P(7)$.}
\StudentItem{From $\exists n\,P(n)$, one may treat the witness as every integer.}
\StudentItem{$c$ is a counterexample to.}
\columnbreak
\StageHeader{3. Expand definitions inside proofs}{Expand the relevant definition before deciding whether a displayed line proves anything.}
\StudentItem{$2k$.}
\StudentItem{$9=2(4)+1$ proves that $9$ is odd.}
\StudentItem{$\sqrt2$ is rational.}
\StudentItem{$4\mid$.}
\StudentItem{The sum of two even integers is even.}
\StageHeader{4. Audit proof language and logical structure}{Separate established facts from goals, conjectures, and incomplete proof language.}
\StudentItem{The product of an even integer and an odd integer is odd.}
\StudentItem{theorem $\Rightarrow$.}
\StudentItem{In $P\to Q$, $Q$ is sufficient for $P$.}
\StudentItem{Using the desired conclusion as an unproved premise is circular.}
\StudentItem{$\forall n\in\mathbb Z,\ n^2\ge0$.}
}
\CoreForm{Proofs}{C}{Variables range over $\mathbb Z$ when bound. $P$ and $Q$ denote propositions, and $P(n)$ denotes a proposition about an integer $n$. A formula with an unbound variable is not a complete proposition.}{$\forall$: for every; $\exists$: there exists; $P\to Q$: if $P$, then $Q$; $P\leftrightarrow Q$: both directions; a \emph{witness} proves an existential; a \emph{counterexample} disproves a universal.}{
\StageHeader{1. Identify the proof strategy}{State what may be assumed and what must be shown for the named strategy.}
\StudentItem{Assume $P$; therefore.}
\StudentItem{Proving $Q\to P$ proves $P\to Q$.}
\StudentItem{$\neg P\Rightarrow$.}
\StudentItem{A proof by cases must cover every possibility allowed by the hypotheses.}
\StudentItem{To prove $P\leftrightarrow Q$, it is enough to prove only $P\to Q$.}
\StageHeader{2. Read quantifiers, witnesses, and counterexamples}{Match universal claims with arbitrary objects and existential claims with witnesses.}
\StudentItem{A proof of a universal statement may begin with an arbitrary object in the domain.}
\StudentItem{To prove $\exists n\,P(n)$, one valid witness is sufficient.}
\StudentItem{From $\exists n\,P(n)$, one may infer $P(7)$ without further information.}
\StudentItem{Let $c$ witness $\exists n$.}
\StudentItem{One integer $c$ with $\neg P(c)$ disproves $\forall n\,P(n)$.}
\columnbreak
\StageHeader{3. Expand definitions inside proofs}{Expand the relevant definition before deciding whether a displayed line proves anything.}
\StudentItem{$6=2(3)$ proves that $6$ is even.}
\StudentItem{$6=2(3)+1$ proves that $6$ is odd.}
\StudentItem{$\frac{3}{0}$ is rational.}
\StudentItem{$4\mid20$.}
\StudentItem{The sum of two odd integers is odd.}
\StageHeader{4. Audit proof language and logical structure}{Separate established facts from goals, conjectures, and incomplete proof language.}
\StudentItem{odd $\cdot$.}
\StudentItem{A theorem is a statement established by proof.}
\StudentItem{In $P\to$, the consequent is necessary.}
\StudentItem{A proof is complete when its final claim merely repeats an unsupported first claim.}
\StudentItem{$\forall n\in\mathbb Z,\ n^2<0$.}
}
\CoreForm{Divisibility}{A}{Let $a,b,c$ be arbitrary integers and let $m\ge2$ be a fixed integer. All other variables are integers when quantified or assigned. In $a\mid b$, the left integer divides the right integer. In $a\equiv_m b$, the modulus satisfies $m\ge2$. A base-$b$ numeral may use only digits $0,\ldots,b-1$.}{$a\mid b$: $a$ divides $b$; $a\nmid b$: $a$ does not divide $b$; $a\equiv b\pmod m$: $a$ and $b$ have the same remainder modulo $m$; $\gcd(a,b)$: greatest common divisor; $(d_k\cdots d_0)_b$: a base-$b$ numeral.}{
\StageHeader{1. Read divisibility from its definition}{Read $a\mid b$ as an existential equation: $b=ak$ for some integer $k$.}
\StudentItem{$3\mid12$.}
\StudentItem{$4\mid20\iff\exists k\in\mathbb Z,\ 20=4k$.}
\StudentItem{$0\mid$.}
\StudentItem{$a\mid b\mid$.}
\StudentItem{If $a\mid b$, then $a\mid(b+1)$.}
\StageHeader{2. Read gcd, quotients, and remainders}{Identify the divisor, quotient, remainder, and required bounds before judging.}
\StudentItem{$\gcd(12,18)=6$.}
\StudentItem{$\gcd(,15)=1$.}
\StudentItem{$23=5(3)+8$ with $0\le8<5$.}
\StudentItem{$23\bmod$.}
\StudentItem{$17\equiv_5 2$.}
\columnbreak
\StageHeader{3. Interpret modular congruence}{Translate congruence into divisibility of a difference before using a law.}
\StudentItem{$a\equiv_m b\iff a\mid(m-b)$.}
\StudentItem{$a\equiv_m b\Rightarrow a+c\equiv_m b$.}
\StudentItem{$a\equiv_m b\Rightarrow ac\equiv_m bc$.}
\StudentItem{If $ac\equiv_m bc$, then $a\equiv_m b$ without any condition on $c$.}
\StudentItem{$14=2k+1$.}
\StageHeader{4. Read primes, bases, and place value}{Check domains, legal digits, and place values before performing arithmetic.}
\StudentItem{$13$ is prime.}
\StudentItem{$(102)_2=6$.}
\StudentItem{$(210)_3=19$.}
\StudentItem{$202=2()+0$.}
\StudentItem{$(10)_b(10)_b=(1000)_b$ for every base $b\ge2$.}
}
\CoreForm{Divisibility}{B}{Let $a,b,c$ be arbitrary integers and let $m\ge2$ be a fixed integer. All other variables are integers when quantified or assigned. In $a\mid b$, the left integer divides the right integer. In $a\equiv_m b$, the modulus satisfies $m\ge2$. A base-$b$ numeral may use only digits $0,\ldots,b-1$.}{$a\mid b$: $a$ divides $b$; $a\nmid b$: $a$ does not divide $b$; $a\equiv b\pmod m$: $a$ and $b$ have the same remainder modulo $m$; $\gcd(a,b)$: greatest common divisor; $(d_k\cdots d_0)_b$: a base-$b$ numeral.}{
\StageHeader{1. Read divisibility from its definition}{Read $a\mid b$ as an existential equation: $b=ak$ for some integer $k$.}
\StudentItem{$5\mid12$.}
\StudentItem{$6\mid20\iff\exists k\in\mathbb Z,\ 20=5k$.}
\StudentItem{$7\mid0$.}
\StudentItem{If $a\mid b$ and $b\mid c$, then $a\mid c$.}
\StudentItem{$a\mid(3b-)$.}
\StageHeader{2. Read gcd, quotients, and remainders}{Identify the divisor, quotient, remainder, and required bounds before judging.}
\StudentItem{$\gcd(12,18)=12$.}
\StudentItem{$\gcd(8,15)=1$.}
\StudentItem{$23=5q+r$.}
\StudentItem{$23\bmod 5=3$.}
\StudentItem{$17\equiv_5 3$.}
\columnbreak
\StageHeader{3. Interpret modular congruence}{Translate congruence into divisibility of a difference before using a law.}
\StudentItem{$a\equiv b$.}
\StudentItem{$a+c\equiv_m$.}
\StudentItem{$a\equiv_m b\Rightarrow ac\equiv_m b$.}
\StudentItem{$\gcd(c,m)=\Rightarrow a\equiv_m b$.}
\StudentItem{$14$ is even.}
\StageHeader{4. Read primes, bases, and place value}{Check domains, legal digits, and place values before performing arithmetic.}
\StudentItem{$15$ is prime.}
\StudentItem{$(101)_2=5$.}
\StudentItem{$(31)_3=10$.}
\StudentItem{$202=2(101)+0$ is a valid first binary-conversion step.}
\StudentItem{$(10)_1(10)_1=(100)_1$.}
}
\CoreForm{Divisibility}{C}{Let $a,b,c$ be arbitrary integers and let $m\ge2$ be a fixed integer. All other variables are integers when quantified or assigned. In $a\mid b$, the left integer divides the right integer. In $a\equiv_m b$, the modulus satisfies $m\ge2$. A base-$b$ numeral may use only digits $0,\ldots,b-1$.}{$a\mid b$: $a$ divides $b$; $a\nmid b$: $a$ does not divide $b$; $a\equiv b\pmod m$: $a$ and $b$ have the same remainder modulo $m$; $\gcd(a,b)$: greatest common divisor; $(d_k\cdots d_0)_b$: a base-$b$ numeral.}{
\StageHeader{1. Read divisibility from its definition}{Read $a\mid b$ as an existential equation: $b=ak$ for some integer $k$.}
\StudentItem{$3\mid$.}
\StudentItem{$20=4k$.}
\StudentItem{$0\mid7$.}
\StudentItem{If $a\mid b$, then $b\mid a$.}
\StudentItem{If $a\mid b$ and $a\mid c$, then $a\mid(3b-2c)$.}
\StageHeader{2. Read gcd, quotients, and remainders}{Identify the divisor, quotient, remainder, and required bounds before judging.}
\StudentItem{$\gcd(12,)$.}
\StudentItem{$\gcd(8,12)=1$.}
\StudentItem{$23=5(4)+3$ with $0\le3<5$.}
\StudentItem{$23\bmod 5=5$.}
\StudentItem{$17\equiv_5$.}
\columnbreak
\StageHeader{3. Interpret modular congruence}{Translate congruence into divisibility of a difference before using a law.}
\StudentItem{$a\equiv_m b\iff m\mid(a-b)$.}
\StudentItem{$a\equiv_m b\Rightarrow a+c\equiv_m b+c$.}
\StudentItem{$ac\equiv_m bc\land$.}
\StudentItem{If $ac\equiv_m bc$ and $\gcd(c,m)=1$, then $a\equiv_m b$.}
\StudentItem{$14$ is odd.}
\StageHeader{4. Read primes, bases, and place value}{Check domains, legal digits, and place values before performing arithmetic.}
\StudentItem{$1$ is prime because.}
\StudentItem{$(101)_2=6$.}
\StudentItem{$(210)_3=21$.}
\StudentItem{$202=2(100)+2$ is a valid first binary-conversion step.}
\StudentItem{$(10)_b(10)_b=(100)_b$ for every base $b\ge2$.}
}
\CoreForm{Induction}{A}{$P(n)$ is a proposition indexed by integers $n\ge1$. In an induction step, $k$ is a fixed but arbitrary integer in the stated range. Also let $a_1=2$ and $a_{n+1}=a_n+3$ for $n\ge1$.}{$P(k)$: the statement at index $k$; $P(k)\to P(k+1)$: the induction step; $a_n$: the $n$th sequence term; $\sum$: add the indexed terms; strong induction may use all earlier established cases.}{
\StageHeader{1. Read the parts of an induction proof}{Label the base case, hypothesis, and step; keep assumptions separate from goals.}
\StudentItem{For a theorem beginning at $n=1$, the base case is $P(1)$.}
\StudentItem{The induction hypothesis assumes $P(k+1)$ in order to prove $P(k+1)$.}
\StudentItem{The induction step proves only $P(k)\to P(k)$.}
\StudentItem{$P(1)$ and $\forall k\ge1\,[P(k)\to P(k+1)]$ imply $\forall n\ge1\,P(n)$.}
\StudentItem{Checking several initial cases does not prove every later case.}
\StageHeader{2. Track hypotheses and base coverage}{Track exactly which earlier cases are available and which starting cases are required.}
\StudentItem{Let $k=$ be arbitrary.}
\StudentItem{Strong induction assumes $P(k+1)$ to prove $P(k+1)$.}
\StudentItem{$P(1)\land\cdots\land P(k)\to$.}
\StudentItem{A recurrence using two preceding terms may require two initial cases.}
\StudentItem{$a_{n+1}=a_n+$.}
\columnbreak
\StageHeader{3. Read recurrences and closed forms}{Substitute an index explicitly before comparing a recurrence with a closed form.}
\StudentItem{$a_1=$.}
\StudentItem{$a_4=11$.}
\StudentItem{$a_n=3n-1$ for every $n\ge1$.}
\StudentItem{The induction hypothesis may replace a different expression that it does not mention.}
\StudentItem{$1+\cdots+n=\frac{n(n+1)}2$.}
\StageHeader{4. Audit the algebra and proof logic}{Check the claimed range and verify that each algebraic line follows from the hypothesis.}
\StudentItem{$2^n<n$ for every integer $n\ge1$.}
\StudentItem{$P(k),\ P(k)\to\ \therefore P(k+1)$.}
\StudentItem{Assuming $P(k+1)$ is allowed while proving $P(k+1)$.}
\StudentItem{A theorem stated for $n\ge4$ is proved by checking only $n=1$.}
\StudentItem{$P(k)+1$.}
}
\CoreForm{Induction}{B}{$P(n)$ is a proposition indexed by integers $n\ge1$. In an induction step, $k$ is a fixed but arbitrary integer in the stated range. Also let $a_1=2$ and $a_{n+1}=a_n+3$ for $n\ge1$.}{$P(k)$: the statement at index $k$; $P(k)\to P(k+1)$: the induction step; $a_n$: the $n$th sequence term; $\sum$: add the indexed terms; strong induction may use all earlier established cases.}{
\StageHeader{1. Read the parts of an induction proof}{Label the base case, hypothesis, and step; keep assumptions separate from goals.}
\StudentItem{For a theorem beginning at $n=1$, the base case is necessarily $P(0)$.}
\StudentItem{Assume $P(k)\land$.}
\StudentItem{$P(k)\to P()$.}
\StudentItem{$P(1)$ alone implies $\forall n\ge1\,P(n)$.}
\StudentItem{Checking $P(1),P(2),P(3)$ proves $\forall n\ge1\,P(n)$.}
\StageHeader{2. Track hypotheses and base coverage}{Track exactly which earlier cases are available and which starting cases are required.}
\StudentItem{The induction step must work for an arbitrary allowed value of $k$.}
\StudentItem{Assume $P(1),\ldots,$.}
\StudentItem{Ordinary and strong induction can prove the same class of statements.}
\StudentItem{Every induction proof has exactly one base case.}
\StudentItem{A recurrence defines later terms using earlier terms.}
\columnbreak
\StageHeader{3. Read recurrences and closed forms}{Substitute an index explicitly before comparing a recurrence with a closed form.}
\StudentItem{The recurrence $a_{n+1}=a_n+3$ and $a_1=2$ determine a unique sequence.}
\StudentItem{$a_4=8$.}
\StudentItem{$a_n=3n+1$ for every $n\ge1$.}
\StudentItem{$1+\cdots+k=$.}
\StudentItem{$\forall n\ge1,\ 1+\cdots+n=\frac{n(n+1)}2$.}
\StageHeader{4. Audit the algebra and proof logic}{Check the claimed range and verify that each algebraic line follows from the hypothesis.}
\StudentItem{$2^n\ge$.}
\StudentItem{From $P(k)$ and $P(k)\to P(k+1)$, one may infer $P(k+1)$.}
\StudentItem{Assume $P(k+1)$ because.}
\StudentItem{$n\ge$.}
\StudentItem{$P(1)\lor\neg P(1)$.}
}
\CoreForm{Induction}{C}{$P(n)$ is a proposition indexed by integers $n\ge1$. In an induction step, $k$ is a fixed but arbitrary integer in the stated range. Also let $a_1=2$ and $a_{n+1}=a_n+3$ for $n\ge1$.}{$P(k)$: the statement at index $k$; $P(k)\to P(k+1)$: the induction step; $a_n$: the $n$th sequence term; $\sum$: add the indexed terms; strong induction may use all earlier established cases.}{
\StageHeader{1. Read the parts of an induction proof}{Label the base case, hypothesis, and step; keep assumptions separate from goals.}
\StudentItem{The base case is $P()$.}
\StudentItem{The induction hypothesis assumes $P(k)$ for an arbitrary $k\ge1$.}
\StudentItem{The induction step proves $P(k)\to P(k+1)$.}
\StudentItem{$P(1)\land\forall k\ge1\Rightarrow\forall n\,P(n)$.}
\StudentItem{$P(1),P(2),P(3)\therefore$.}
\StageHeader{2. Track hypotheses and base coverage}{Track exactly which earlier cases are available and which starting cases are required.}
\StudentItem{Proving $P(7)\to P(8)$ establishes the complete induction step for every $k$.}
\StudentItem{Strong induction may assume $P(1),\ldots,P(k)$ to prove $P(k+1)$.}
\StudentItem{Strong induction is valid only when ordinary induction fails.}
\StudentItem{Base cases: $P(1),$.}
\StudentItem{A recurrence alone always determines a unique sequence without initial values.}
\columnbreak
\StageHeader{3. Read recurrences and closed forms}{Substitute an index explicitly before comparing a recurrence with a closed form.}
\StudentItem{The recurrence $a_{n+1}=a_n+3$ determines one unique sequence without $a_1$.}
\StudentItem{$a_{4}+$.}
\StudentItem{$a_n=3n-1$.}
\StudentItem{In a sum proof, the induction hypothesis may replace the partial sum through $k$.}
\StudentItem{$\forall n\ge1,\ 1+\cdots+n=\frac{n(n-1)}2$.}
\StageHeader{4. Audit the algebra and proof logic}{Check the claimed range and verify that each algebraic line follows from the hypothesis.}
\StudentItem{$2^n\ge n+1$ for every integer $n\ge0$.}
\StudentItem{From $P(k+1)$ and $P(k)\to P(k+1)$, one may infer $P(k)$.}
\StudentItem{Assuming $P(k)$ is allowed while proving $P(k)\to P(k+1)$.}
\StudentItem{A theorem stated for $n\ge4$ should begin its induction at $n=4$.}
\StudentItem{$P(1)\land\neg P(1)$.}
}
\CoreForm{Counting}{A}{Let $X=\{a,b\}$, $Y=\{b,c,d\}$, $n=5$, and $r=2$. For probability, a fair die has sample space $\Omega=\{1,2,3,4,5,6\}$ with events $E=\{2,4,6\}$ and $F=\{4,5,6\}$. Let $Z$ be the number shown.}{$n!$: factorial; $P(n,r)$: ordered selections; $\binom nr$: unordered selections; $|X|$: size of $X$; $\Pr(E)$: probability of event $E$; $E^c$: not $E$; $\Pr(E\mid F)$: probability of $E$ given $F$; $\mathbb E[Z]$: expected value.}{
\StageHeader{1. Choose and read a counting rule}{Decide whether stages are combined by addition or multiplication before calculating.}
\StudentItem{$3+=$.}
\StudentItem{A two-stage process with $3$ choices and then $4$ choices has $7$ outcomes.}
\StudentItem{$|X\cup Y|=5$.}
\StudentItem{$7\to6\therefore$.}
\StudentItem{$X!$.}
\StageHeader{2. Interpret counting notation and models}{Ask whether order and repetition matter, then match the notation to that model.}
\StudentItem{$P(5,2)=20$.}
\StudentItem{$\binom52=20$.}
\StudentItem{$\binom52=\binom53$.}
\StudentItem{There are $2^5=32$ bit strings of length $5$.}
\StudentItem{Five choices made twice with repetition produce $P(5,2)$ sequences.}
\columnbreak
\StageHeader{3. Read events and probability rules}{Identify the sample space and event type before applying a probability rule.}
\StudentItem{$\frac{5!}{2!\,}$.}
\StudentItem{Selecting a two-person committee from five people is counted by $\binom52$.}
\StudentItem{$7\in\Omega$.}
\StudentItem{$\Pr(E)=\frac13$.}
\StudentItem{$\Pr(E^c)=1-\Pr(E)$.}
\StageHeader{4. Read dependence, conditioning, and expectation}{Check conditioning, dependence, and the type of the object inside each operator.}
\StudentItem{$\Pr(E\cup)=1$.}
\StudentItem{$E\cap$ is empty.}
\StudentItem{$\Pr(E\cap F)=\frac13$.}
\StudentItem{$\Pr(E\mid\varnothing)$.}
\StudentItem{$\mathbb E[Z]=3$.}
}
\CoreForm{Counting}{B}{Let $X=\{a,b\}$, $Y=\{b,c,d\}$, $n=5$, and $r=2$. For probability, a fair die has sample space $\Omega=\{1,2,3,4,5,6\}$ with events $E=\{2,4,6\}$ and $F=\{4,5,6\}$. Let $Z$ be the number shown.}{$n!$: factorial; $P(n,r)$: ordered selections; $\binom nr$: unordered selections; $|X|$: size of $X$; $\Pr(E)$: probability of event $E$; $E^c$: not $E$; $\Pr(E\mid F)$: probability of $E$ given $F$; $\mathbb E[Z]$: expected value.}{
\StageHeader{1. Choose and read a counting rule}{Decide whether stages are combined by addition or multiplication before calculating.}
\StudentItem{If two disjoint cases have $3$ and $4$ outcomes, together they have $7$ outcomes.}
\StudentItem{$3\cdot$.}
\StudentItem{$|X\cup|=4$.}
\StudentItem{Placing $7$ objects into $6$ boxes forces some box to receive at least two objects.}
\StudentItem{$5!=120$.}
\StageHeader{2. Interpret counting notation and models}{Ask whether order and repetition matter, then match the notation to that model.}
\StudentItem{$P(5,2)=10$.}
\StudentItem{$\binom{\,}{2}=10$.}
\StudentItem{$\binom52=\binom54$.}
\StudentItem{There are $5!=120$ bit strings of length $5$.}
\StudentItem{$5^{\Box}$ sequences.}
\columnbreak
\StageHeader{3. Read events and probability rules}{Identify the sample space and event type before applying a probability rule.}
\StudentItem{The word LEVEL has $\frac{5!}{2!\,2!}=30$ distinct arrangements.}
\StudentItem{Assigning president and secretary from five people is counted by $\binom52$.}
\StudentItem{$E\in 4$.}
\StudentItem{$\Pr(2)=\frac16$.}
\StudentItem{$\Pr(E^c)=\Pr(E)+1$.}
\StageHeader{4. Read dependence, conditioning, and expectation}{Check conditioning, dependence, and the type of the object inside each operator.}
\StudentItem{$\Pr(E\cup F)=\Pr(E)+\Pr(F)-\Pr(E\cap F)$.}
\StudentItem{$\{1\}$ and $\{2\}$ are disjoint events.}
\StudentItem{$E$ and $F$ are independent.}
\StudentItem{$\Pr(E\mid F)=\frac23$.}
\StudentItem{$\mathbb E[E]$.}
}
\CoreForm{Counting}{C}{Let $X=\{a,b\}$, $Y=\{b,c,d\}$, $n=5$, and $r=2$. For probability, a fair die has sample space $\Omega=\{1,2,3,4,5,6\}$ with events $E=\{2,4,6\}$ and $F=\{4,5,6\}$. Let $Z$ be the number shown.}{$n!$: factorial; $P(n,r)$: ordered selections; $\binom nr$: unordered selections; $|X|$: size of $X$; $\Pr(E)$: probability of event $E$; $E^c$: not $E$; $\Pr(E\mid F)$: probability of $E$ given $F$; $\mathbb E[Z]$: expected value.}{
\StageHeader{1. Choose and read a counting rule}{Decide whether stages are combined by addition or multiplication before calculating.}
\StudentItem{If two disjoint cases have $3$ and $4$ outcomes, together they have $12$ outcomes.}
\StudentItem{A two-stage process with $3$ choices and then $4$ choices has $12$ outcomes.}
\StudentItem{$|X\cup Y|=4$.}
\StudentItem{Placing $6$ objects into $6$ boxes forces some box to receive at least two objects.}
\StudentItem{$5!=25$.}
\StageHeader{2. Interpret counting notation and models}{Ask whether order and repetition matter, then match the notation to that model.}
\StudentItem{$P(5,)=20$.}
\StudentItem{$\binom52=10$.}
\StudentItem{$\binom5X=10$.}
\StudentItem{$2^{\Box}=32$.}
\StudentItem{Five choices made twice with repetition produce $5^2=25$ sequences.}
\columnbreak
\StageHeader{3. Read events and probability rules}{Identify the sample space and event type before applying a probability rule.}
\StudentItem{The word LEVEL has $5!=120$ distinct arrangements.}
\StudentItem{committee$(5,)$.}
\StudentItem{$E\subseteq\Omega$.}
\StudentItem{$\Pr(E)=\frac12$.}
\StudentItem{$\Pr({}^c)=1$.}
\StageHeader{4. Read dependence, conditioning, and expectation}{Check conditioning, dependence, and the type of the object inside each operator.}
\StudentItem{$\Pr(E\cup F)=\Pr(E)+\Pr(F)$.}
\StudentItem{$E$ and $F$ are disjoint events.}
\StudentItem{$\Pr(E)\Pr()=\Pr(E\cap F)$.}
\StudentItem{$\Pr(E\mid F)=\frac12$.}
\StudentItem{$\mathbb E[Z]=3.5$.}
}

\end{document}
